\documentclass[10pt,a4paper]{amsart}
\usepackage[margin=19mm]{geometry}
\usepackage{amsmath,amssymb,mathtools}
\usepackage[scr=rsfs,cal=euler]{mathalfa}
\usepackage{xfrac}
\usepackage[unicode,hidelinks,bookmarks=false]{hyperref}
\newcommand{\ie}{i.e.\ }
\newcommand{\cf}{cf.\ }
\newcommand{\resp}{respectively\ }
\newcommand{\Subsection}{\subsection}
\providecommand{\nobreakdash}{\nobreak\hbox{-}}
\setlength{\emergencystretch}{2em}
\multlinegap0pt
\usepackage{tikz}
\usetikzlibrary{cd}
\usepackage[normalem]{ulem}
\usepackage{mathtools}
\usepackage{bm}
\usepackage{tensor}
\usepackage{enumerate}
\usepackage{braket}
\usepackage{csquotes}
\usepackage[shortlabels]{enumitem}
\usepackage{tikz}
\usepackage{xcolor}
\newtheorem{lemma}{Lemma}
\newtheorem{proposition}{Proposition}
\usepackage{cleveref}
\crefname{lemma}{Lemma}{Lemmas}
\crefname{proposition}{Proposition}{Propositions}
\newcommand{\LPr}{\mathrm{Pr}^\rL}
\newcommand{\Mod}{\mathrm{Mod}}
\newcommand{\Nuc}{\mathrm{Nuc}}
\newcommand{\Qcoh}{\mathrm{QCoh}}
\newcommand{\Spec}{\mathrm{Spec}}
\newcommand{\bDelta}{\bm{\Delta}} % Required for inserting images
\newcommand{\cat}{$\infty$-category }
 \newcommand{\rL}{\mathrm L}
\title{Smooth categories in a 6 functor formalism and compact generation for nuclear categories in analytic geometry}
\author{Matteo Montagnani}
\date{Source: arXiv:2605.21024v1 (CC BY 4.0)}
\begin{document}
\maketitle

\noindent\emph{Source and licence.} Smooth categories in a 6 functor formalism and compact generation for nuclear categories in analytic geometry. Version arXiv:2605.21024v1 (CC BY 4.0), distributed by arXiv under the Creative Commons Attribution 4.0 International licence, http://creativecommons.org/licenses/by/4.0/, as recorded in the arXiv licence metadata for this exact version and shown on the abstract page https://arxiv.org/abs/2605.21024v1. Full licence notice: \texttt{licenses/arxiv-2605.21024-CC-BY-4.0.txt}.\par\smallskip

\noindent\textit{Editorial scope.} Counted unit: the article's proposition GAGA (source0.tex lines 2409--2411). The article's complete original proof of it is delivered with it (source0.tex lines 2412--2428, one \texttt{proof} environment of 17 lines). No mathematics is written, repaired or re-derived: the statement and the proof are the article's own lines, in the article's own notation, and no retained line is substituted or re-flowed.  Retained uncounted as the source-local context the proof needs: lines 2383--2393.  Nothing was omitted from any retained span, so the package carries no omission record.  No retained line contains a citation, so the wrapper carries no bibliography and no bibliography entry.  The retained lines contain the article's own diagram code, which the wrapper typesets with the standard tikz packages; no image file is delivered and the package declares no asset dependency.  Editorial formatting only, in addition: an AMS \texttt{amsart} preamble that restates the article's own declarations and macros used by the retained lines, namely the environments \texttt{lemma}, \texttt{proposition} on separate counters, together with the standard packages those lines need.  The retained environments print in this document as Lemma 1, Proposition 1 in order of appearance, whereas the article prints the same items with the numbering of its own full document.  The complete native source of the article as distributed in the arXiv e-print for this exact version is bundled unchanged as \texttt{sources/201035/source0.tex}.
\begin{lemma}\label{lem: descent for discrete modules}
    The functor
    \begin{equation}\label{eq: discrete descent}
        \begin{tikzcd}
	{ \mathrm{CAlg}(\Mod_{k})} && {\mathrm{AnRing}} && \LPr
	\arrow["{(-)^{\mathrm{k-alg}}}", from=1-1, to=1-3]
	\arrow["{\Nuc(-)}", from=1-3, to=1-5]
\end{tikzcd}
    \end{equation}
    defined by sending a commutative algebra $R$ to the \cat of nuclear modules over the analytic ring $R^{\mathrm{k-alg}}$, satisfies descent for the fppf topology.
\end{lemma}

\begin{proposition}\label{prop: algebraic GAGA}
    Let $X$ be a connective derived scheme, separated and of finite type over $k$. Then there is an equivalence of categories between $\Nuc(X^{\mathrm{k-alg}})$ and $\Qcoh(X)\otimes_{\Mod_{k}}\Nuc(k)$.
\end{proposition}

\begin{proof}
    As above, let $\{\Spec(A_{i})\}_{i \in I}$ be a cover of $X$, and let $\Spec(A_{i_1, \dots, i_{n}})$ denote the intersection $\Spec(A_{i_{1}})\times_{X}\dots \times_{X}\Spec(A_{i_{n}})$. Note that since $X$ is separated, $\Qcoh(X)$ can be written as a limit of quasi-coherent sheaves on affine schemes. In particular, using the same notation as above, we have that 
    \[
    \Qcoh(X) \simeq \lim_{[n]\in\bDelta^{\mathrm{op}}}\prod_{\left\{i_1,\cdots,i_n\right\}\subseteq I}\Mod_{A_{i_1, \dots, i_{n}}}.
    \]
    Thus, $\Qcoh(X)$ can be written as a limit of $\Mod_{A}$ for affine opens $\Spec(A)$ in $X$.
    Now, since both $\Nuc(k)$ and $\Mod_{k}$ are rigid, it follows that $\Nuc(k)$ is dualizable over $\Mod_{k}$; in particular, the functor 
    $-\otimes_{\Mod_{k}}\Nuc(k)$ commutes with limits. Therefore, we can write 
    \begin{equation}\label{eq:Kunneth limit}
        \Qcoh(X)\otimes_{\Mod_{k}}\Nuc(k) \simeq \lim_{[n]\in\bDelta^{\mathrm{op}}}\prod_{\left\{i_1,\cdots,i_n\right\}\subseteq I}\left(\Mod_{A_{i_1, \dots, i_{n}}}\otimes_{\Mod_k}\Nuc(k)\right).
    \end{equation}
    The right-hand side can be identified with 
    \[
     \lim_{[n]\in\bDelta^{\mathrm{op}}}\prod_{\left\{i_1,\cdots,i_n\right\}\subseteq I}\Mod_{\uline{A}_{i_1, \dots, i_{n}}}(\Nuc(k)).
    \]
    To conclude, it is enough to observe that $\Mod_{\uline{A}}(\Nuc(k))$ can be identified with $\Nuc(A^{\mathrm{k-alg}})$, and that $\{\mathrm{AnSpec}(A_i^{\mathrm{k-alg}})\}_{i \in I}$ defines a cover of $X^{\mathrm{k-alg}}$. Hence, using descent for $\Nuc(-)$ (\Cref{lem: descent for discrete modules}), we can identify the limit in \Cref{eq:Kunneth limit} with $\Nuc(X^{\mathrm{k-alg}})$.
\end{proof}

\end{document}
